%% =========================================================
%% Chapter: The Algebraic Fundamental Theorem
%% Covers the fixed-system-size approach of arXiv:1804.04964
%% (Moln\'ar, Garre-Rubio, P\'erez-Garc\'ia, Schuch, Cirac, NJP 2018)
%% Stages 1-5c of tracking issue #11.
%% Chain definitions (Stage 0) live in ch03_single.tex.
%% =========================================================

\chapter{The Algebraic Fundamental Theorem}%
\label{ch:algebraic_ft}

This chapter develops the algebraic proof of the Fundamental Theorem for
injective MPS on a finite periodic chain of fixed length, following
\cite{Molnar2018NormalPEPS}.
The approach is self-contained at the level of linear algebra: it never
appeals to operator-norm limits or spectral theory of transfer maps,
working instead with the matrix algebra structure directly.

The chain bookkeeping layer (Definition~\ref{def:non_ti_chain}),
including the notions of same chain state, injectivity, and cyclic gauge
equivalence, was introduced in Chapter~\ref{ch:single}.
This chapter takes those definitions as given and proceeds through five
stages:
\begin{enumerate}
    \item construction of a one-sided inverse of the linear combination map
          (Section~\ref{sec:decomposition_map}),
    \item physical realisation of virtual bond insertions
          (Section~\ref{sec:phys_realize}),
    \item proof that the same state forces bond algebras to be isomorphic
          (Section~\ref{sec:virtual_bond_gauge}),
    \item proof that aligned gauges force proportionality of local tensors
          (Section~\ref{sec:tensor_proportional}),
    \item assembly into the Fundamental Theorem and corollaries
          (Sections~\ref{sec:ft_chain} and~\ref{sec:ft_corollaries}).
\end{enumerate}

%% ---------------------------------------------------------
\section{One-sided inverse and decomposition}%
\label{sec:decomposition_map}
%% ---------------------------------------------------------

Let $A = \{A^i\}_{i=0}^{d-1}$ be an injective MPS tensor with bond
dimension $D$.
Recall from Definition~\ref{def:injective} that injectivity means that
the matrices $\{A^i\}$ span $\MN{D}$.
In other words, the \emph{linear combination map}
\[
    \Phi_A : \C^d \to \MN{D},
    \qquad
    \Phi_A(c) := \sum_{i=0}^{d-1} c_i A^i,
\]
is surjective.
A right inverse $\Psi_A$ of $\Phi_A$ therefore exists, and we record it
as a Lean definition.

\begin{definition}[Decomposition map]\label{def:decomposition_map}
    \lean{MPSTensor.decompositionMap}
    \leanok
    \uses{def:injective}
    Let $A$ be an injective MPS tensor.
    The \emph{decomposition map} is a fixed right inverse
    \[
        \Psi_A : \MN{D} \to \C^d
    \]
    of the linear combination map $\Phi_A$, i.e.\
    $\Phi_A \circ \Psi_A = \mathrm{id}_{\MN{D}}$.
    Such a map exists because $\Phi_A$ is surjective.
\end{definition}

\begin{lemma}[Decomposition property]\label{lem:exists_decomposition}
    \lean{MPSTensor.IsInjective.exists_decomposition}
    \leanok
    \uses{def:decomposition_map}
    If $A$ is injective, then for every $X \in \MN{D}$ there exist
    coefficients $c \in \C^d$ such that
    \[
        X = \sum_{i=0}^{d-1} c_i\, A^i.
    \]
\end{lemma}

\begin{proof}
    \leanok
    Apply the decomposition map: take $c = \Psi_A(X)$.
    Then $\Phi_A(c) = \Phi_A(\Psi_A(X)) = X$.
\end{proof}

\begin{remark}[Non-uniqueness of the decomposition]
    \label{rem:decomposition_nonunique}
    When $d > D^2$ the map $\Phi_A$ has a nontrivial kernel and the
    decomposition is not unique.
    The decomposition map $\Psi_A$ is any fixed section; its choice
    does not affect the statements that follow because those statements
    only depend on the image $\Phi_A \circ \Psi_A = \mathrm{id}$.
    Uniqueness of the decomposition holds if and only if $d = D^2$ and
    the matrices $\{A^i\}$ form a basis of $\MN{D}$.
\end{remark}

%% ---------------------------------------------------------
\section{Physical realisation of virtual insertions}%
\label{sec:phys_realize}
%% ---------------------------------------------------------

Given an injective tensor $A$ and a matrix $X \in \MN{D}$ to be
inserted on the virtual bond, one can ask: is there a $d \times d$ matrix
$O$ such that inserting $X$ between sites $k$ and $k+1$ on the virtual
bond has the same effect on any multi-site amplitude as acting with $O$
on the physical index at site $k+1$?
The next definition makes this precise and shows the answer is yes.

\begin{definition}[Physical realisation map]\label{def:phys_realize}
    \lean{MPSTensor.physRealize}
    \leanok
    \uses{def:decomposition_map}
    Let $A$ be an injective tensor.
    For $X \in \MN{D}$, define the $d \times d$ matrix
    $O_A(X) \in \MN{d}$ by
    \[
        O_A(X)_{ij} := (\Psi_A(A^i X))_j,
    \]
    so that row~$i$ of $O_A(X)$ contains the decomposition coefficients
    of $A^i X$ in the generating family $\{A^j\}$.
    The map $X \mapsto O_A(X)$ is $\C$-linear in $X$.
\end{definition}

\begin{lemma}[Defining property of the physical realisation map]%
    \label{lem:physRealize_spec}
    \lean{MPSTensor.physRealize_spec}
    \leanok
    \uses{def:phys_realize, def:decomposition_map}
    For every $X \in \MN{D}$ and every physical index $i$,
    \[
        A^i\, X = \sum_{j=0}^{d-1} O_A(X)_{ij}\, A^j.
    \]
\end{lemma}

\begin{proof}
    \leanok
    By definition, $O_A(X)_{ij} = (\Psi_A(A^i X))_j$.
    Applying $\Phi_A$ gives
    $\sum_j O_A(X)_{ij} A^j = \Phi_A(\Psi_A(A^i X)) = A^i X$,
    where the last step uses $\Phi_A \circ \Psi_A = \mathrm{id}$.
\end{proof}

The key algebraic property is that the physical realisation map is a ring
homomorphism with respect to matrix multiplication on the virtual bond.

\begin{theorem}[Physical realisation is multiplicative]%
    \label{thm:physRealize_mul}
    \lean{MPSTensor.physRealize_mul}
    \notready
    \uses{def:phys_realize, lem:physRealize_spec}
    Let $A$ be injective.
    For all $X, Y \in \MN{D}$,
    \[
        O_A(XY) = O_A(X)\, O_A(Y).
    \]
\end{theorem}

\begin{proof}
    By Lemma~\ref{lem:physRealize_spec} applied to $Y$,
    \[
        A^j\, Y = \sum_{k} O_A(Y)_{jk}\, A^k.
    \]
    Apply Lemma~\ref{lem:physRealize_spec} to $XY$ and also to $X$:
    \begin{align*}
        A^i\,(XY)
        &= \sum_j O_A(XY)_{ij}\, A^j, \\
        A^i\, X
        &= \sum_j O_A(X)_{ij}\, A^j.
    \end{align*}
    On the other hand, using associativity and the formula for $A^j Y$,
    \begin{align*}
        A^i\,(XY) = (A^i X)\,Y
        &= \left(\sum_j O_A(X)_{ij}\, A^j\right) Y
         = \sum_j O_A(X)_{ij}\, (A^j Y) \\
        &= \sum_j O_A(X)_{ij} \sum_k O_A(Y)_{jk}\, A^k
         = \sum_k (O_A(X)\, O_A(Y))_{ik}\, A^k.
    \end{align*}
    Thus $\sum_j O_A(XY)_{ij} A^j = \sum_k (O_A(X) O_A(Y))_{ik} A^k$
    for every $i$.
    When the $\{A^j\}$ are linearly independent (i.e.\ $d = D^2$ and
    $\{A^j\}$ form a basis of $\MN{D}$), uniqueness of the decomposition
    gives $O_A(XY) = O_A(X) O_A(Y)$ immediately.
    The general injective case (where $d \ge D^2$ and the decomposition
    may not be unique) requires an additional argument to identify the
    coefficients; this is currently the remaining gap in the Lean
    formalisation.
\end{proof}

%% ---------------------------------------------------------
\section{Same state forces bond algebras to be isomorphic}%
\label{sec:virtual_bond_gauge}
%% ---------------------------------------------------------

This section proves Lemma~1 of \cite{Molnar2018NormalPEPS}: if two
injective MPS on a chain of length $n \ge 3$ generate the same state,
then the algebra of virtual bond insertions for the first chain is
isomorphic to that of the second.
The isomorphism is inner, so it is implemented by conjugation by an
invertible bond matrix $Z$.

\begin{theorem}[Virtual bond gauge - Lemma~1 of \protect\cite{Molnar2018NormalPEPS}]%
    \label{thm:virtual_bond_gauge}
    \lean{}
    \notready
    \uses{thm:physRealize_mul, thm:simplicity, thm:skolem_noether,
          def:non_ti_chain}
    Let $\mathbf{A} = (A_0, \dots, A_{n-1})$ and
    $\mathbf{B} = (B_0, \dots, B_{n-1})$ be two injective non-TI
    periodic chains of length $n \ge 3$ that generate the same chain
    state.
    Fix a bond position $k \in \{0, \dots, n-1\}$.
    Then there exists an invertible matrix $Z_k \in \GL_D(\C)$ such that
    for every $X \in \MN{D}$,
    \[
        O_{A, k}(X) = O_{B, k}(Z_k^{-1} X Z_k),
    \]
    where $O_{A,k}$ (resp.\ $O_{B,k}$) denotes the physical realisation
    map obtained by blocking the sites adjacent to bond~$k$ in the
    $\mathbf{A}$-chain (resp.\ $\mathbf{B}$-chain).
    Equivalently, inserting $X$ on bond~$k$ in the $\mathbf{A}$-chain
    gives the same physical operation as inserting $Z_k^{-1} X Z_k$ on
    bond~$k$ in the $\mathbf{B}$-chain.
\end{theorem}

\begin{proof}
    The proof has six steps.

    \textbf{Step 1} (physical realisation for $\mathbf{A}$).
    Block the sites on either side of bond~$k$ into a single three-site
    block.
    Injectivity of each local tensor implies injectivity of the blocked
    tensor (Lemma~\ref{lem:exists_decomposition}).
    Define $O_{A,k}(X) \in \MN{d^2}$ via the physical realisation map
    for this blocked tensor.
    By Theorem~\ref{thm:physRealize_mul}, the map $X \mapsto O_{A,k}(X)$
    is a $\C$-algebra homomorphism from $\MN{D}$ to $\MN{d^2}$.

    \textbf{Step 2} (resonance on the $\mathbf{B}$-chain).
    Since $\mathbf{A}$ and $\mathbf{B}$ generate the same chain state,
    inserting any physical operator on bond~$k$ gives the same amplitude
    for both chains.
    Hence for every $X \in \MN{D}$ the physical action $O_{A,k}(X)$ of
    the $\mathbf{A}$-chain can be realised on the $\mathbf{B}$-chain.
    Formally, there exists a unique $Y \in \MN{D}$ such that
    $O_{A,k}(X) = O_{B,k}(Y)$.

    \textbf{Step 3} (virtual action on $\mathbf{B}$).
    Applying the physical realisation map $O_{B,k}$ to elements of
    $\MN{D}$ and using injectivity of $\mathbf{B}$, the map
    $T : X \mapsto Y$ defined in Step~2 is well-defined and $\C$-linear.

    \textbf{Step 4} (multiplicativity of $T$).
    Since both $O_{A,k}$ and $O_{B,k}$ are algebra homomorphisms
    (Theorem~\ref{thm:physRealize_mul}), the composed map $T$ satisfies
    $T(XY) = T(X)T(Y)$ for all $X, Y \in \MN{D}$.

    \textbf{Step 5} (bijectivity of $T$).
    $T$ is a nonzero multiplicative $\C$-linear endomorphism of $\MN{D}$.
    By Theorem~\ref{thm:simplicity}, $T$ is bijective, hence a
    $\C$-algebra automorphism.

    \textbf{Step 6} (inner conjugate via Skolem--Noether).
    By Theorem~\ref{thm:skolem_noether}, every $\C$-algebra automorphism
    of $\MN{D}$ is inner.
    Hence there exists $Z_k \in \GL_D(\C)$ with
    $T(X) = Z_k^{-1} X Z_k$ for all $X$, which is the claimed relation.
\end{proof}

%% ---------------------------------------------------------
\section{Aligned gauges force proportionality}%
\label{sec:tensor_proportional}
%% ---------------------------------------------------------

After applying Theorem~\ref{thm:virtual_bond_gauge} to each bond, one
can absorb the gauge matrices $Z_k$ into the tensors to produce an
intermediate chain $\widetilde{\mathbf{B}}$ that agrees with $\mathbf{A}$
on all virtual insertions.
Lemma~2 of \cite{Molnar2018NormalPEPS} then shows that two such chains,
whose virtual-insertion operations agree on both bonds of a two-site
block, must be proportional.

\begin{theorem}[Aligned gauges force proportionality - Lemma~2 of \protect\cite{Molnar2018NormalPEPS}]%
    \label{thm:tensor_proportional}
    \lean{}
    \notready
    \uses{def:injective, def:decomposition_map}
    Let $A_0, A_1$ and $B_0, B_1$ each be injective MPS tensors with the
    same bond dimension $D$.
    Suppose that for every matrix $M \in \MN{D}$,
    \[
        A_1^{j'} M A_0^j = B_1^{j'} M B_0^j
        \qquad
        \text{for all physical indices } j, j'.
    \]
    Then there exists a nonzero scalar $\lambda \in \C^{\times}$ such that
    \[
        A_0^j = \lambda\, B_0^j
        \qquad\text{and}\qquad
        A_1^{j'} = \lambda^{-1}\, B_1^{j'}
    \]
    for all physical indices $j, j'$.
\end{theorem}

\begin{proof}
    The proof has five steps.

    \textbf{Step 1} (single bond relation from two-site matching).
    Setting $M = \Id_D$ in the hypothesis gives
    $A_1^{j'} A_0^j = B_1^{j'} B_0^j$ for all $j, j'$.
    In particular, summing over~$j$ with decomposition coefficients,
    Lemma~\ref{lem:exists_decomposition} provides for every
    $X \in \MN{D}$ a vector $c$ with $X = \sum_j c_j B_0^j$, so
    \begin{equation}
        A_1^{j'} \left(\sum_j c_j A_0^j\right) = B_1^{j'} X.
        \label{eq:one_bond}
    \end{equation}

    \textbf{Step 2} (existence of right-hand factor).
    From the general hypothesis $A_1^{j'} M A_0^j = B_1^{j'} M B_0^j$,
    applying the decomposition map $\Psi_{A_0}$ on the $j$-index yields a
    matrix $Z$ with $A_0^j = Z B_0^j$ for all~$j$.
    Similarly, applying $\Psi_{A_1}$ on the $j'$-index yields a matrix
    $W$ with $A_1^{j'} = B_1^{j'} W^{-1}$ for all~$j'$.

    \textbf{Step 3} (commutativity with all matrices).
    Since the $\{B_0^j\}$ span $\MN{D}$ (injectivity of $B_0$), the
    relation $A_0^j = Z B_0^j$ holding for all~$j$ forces $Z$ to commute
    with every matrix in the span, i.e.\ $Z$ commutes with every element
    of $\MN{D}$.

    \textbf{Step 4} (centre of $\MN{D}$ is scalar).
    The centre of $\MN{D}$ consists exactly of scalar multiples of the
    identity: if $ZN = NZ$ for all $N \in \MN{D}$, then $Z = \lambda \Id$
    for some $\lambda \in \C$.

    \textbf{Step 5} (conclusion).
    Hence $Z = \lambda \Id$ and $A_0^j = \lambda B_0^j$ for all~$j$.
    The matching condition with $M = \Id$ then gives
    $A_1^{j'} = \lambda^{-1} B_1^{j'}$ for all~$j'$, as required.
    Since the tensors are nonzero, $\lambda \neq 0$.
\end{proof}

%% ---------------------------------------------------------
\section{The Fundamental Theorem for injective chains}%
\label{sec:ft_chain}
%% ---------------------------------------------------------

Combining the gauge-conjugacy result
(Theorem~\ref{thm:virtual_bond_gauge}) with the proportionality result
(Theorem~\ref{thm:tensor_proportional}) yields the main theorem of
\cite{Molnar2018NormalPEPS} for non-translation-invariant injective MPS.

\begin{theorem}[Fundamental Theorem for injective chains - Theorem~1 of \protect\cite{Molnar2018NormalPEPS}]%
    \label{thm:fundamentalTheorem_injective_chain}
    \lean{}
    \notready
    \uses{thm:virtual_bond_gauge, thm:tensor_proportional,
          def:non_ti_chain}
    Let $\mathbf{A} = (A_0, \dots, A_{n-1})$ and
    $\mathbf{B} = (B_0, \dots, B_{n-1})$ be two injective non-TI periodic
    chains of length $n \ge 3$ that generate the same chain state.
    Then $\mathbf{A}$ and $\mathbf{B}$ are cyclically gauge equivalent:
    there exist invertible matrices $Z_0, \dots, Z_{n-1} \in \GL_D(\C)$
    such that
    \[
        B_k^i = Z_k\, A_k^i\, Z_{k+1}^{-1}
    \]
    for every site $k \in \{0, \dots, n-1\}$ and every physical index
    $i$, with $Z_n = Z_0$.
\end{theorem}

\begin{proof}
    The proof proceeds in four steps.

    \textbf{Step 1} (gauge on each bond via Lemma~1).
    Fix a bond~$k$.
    Block the sites not equal to~$k$ and $k+1$ around the chain to form
    a three-site configuration, which is again injective.
    Apply Theorem~\ref{thm:virtual_bond_gauge} to obtain an invertible
    $Z_k \in \GL_D(\C)$ such that inserting $X$ on bond~$k$ in the
    $\mathbf{A}$-chain is equivalent to inserting $Z_k^{-1} X Z_k$ on
    bond~$k$ in the $\mathbf{B}$-chain.

    \textbf{Step 2} (absorbing gauges).
    Define the modified $\mathbf{B}$-chain by
    $\widetilde{B}_k^i := Z_k^{-1} B_k^i Z_{k+1}$ for each site~$k$.
    After this gauge absorption, the virtual insertion operations for
    $\mathbf{A}$ and $\widetilde{\mathbf{B}}$ agree on every bond.

    \textbf{Step 3} (proportionality via Lemma~2).
    Fix a pair of adjacent sites $(k, k+1)$.
    Block all other sites into a single matrix $M$ running around the
    periodic chain.
    The agreement of virtual insertions on both bonds of this two-site
    block implies the hypothesis of Theorem~\ref{thm:tensor_proportional}.
    Hence there exist scalars $\lambda_k \in \C^{\times}$ with
    $A_k^i = \lambda_k \widetilde{B}_k^i$ for each site~$k$.

    \textbf{Step 4} (absorbing scalars).
    Redefine $Z_k \mapsto \lambda_k Z_k$ (absorbing the scalar into the
    gauge at each bond); the cyclic gauge condition $B_k^i = Z_k A_k^i Z_{k+1}^{-1}$
    is then satisfied for all $k$ and $i$.
\end{proof}

%% ---------------------------------------------------------
\section{Corollaries}%
\label{sec:ft_corollaries}
%% ---------------------------------------------------------

\subsection{Translation-invariant chains}

\begin{corollary}[TI case]\label{cor:ft_ti_chain}
    \lean{}
    \notready
    \uses{thm:fundamentalTheorem_injective_chain, def:non_ti_chain}
    Let $A$ and $B$ be injective MPS tensors of bond dimension $D$.
    Suppose the two translation-invariant chains
    $(A, A, \dots, A)$ and $(B, B, \dots, B)$ of length $n \ge 3$
    generate the same chain state.
    Then there exists an invertible $Z \in \GL_D(\C)$ and a scalar
    $\lambda \in \C^{\times}$ with $\lambda^n = 1$ such that
    \[
        B^i = \lambda\, Z\, A^i\, Z^{-1}
    \]
    for every physical index~$i$.
\end{corollary}

\begin{proof}
    Apply Theorem~\ref{thm:fundamentalTheorem_injective_chain}
    to the two constant chains.
    This yields invertible matrices $Z_0, \dots, Z_{n-1}$ with
    $B^i = Z_k A^i Z_{k+1}^{-1}$ for all $k$ and $i$.
    Setting $G_k = Z_k Z_0^{-1}$, the tensor $A$ satisfies
    $A^i = G_k^{-1} A^i G_k$ for all $k$, so each $G_k$ commutes
    with all elements of $\MN{D}$ (since the $\{A^i\}$ span $\MN{D}$).
    Hence $G_k = \mu_k \Id$ for some scalar $\mu_k$, giving
    $Z_k = \mu_k Z_0$.
    The cyclic condition $Z_n = Z_0$ forces $\mu_n = 1$, and the ratio
    $\mu_1 = \mu_k / \mu_{k-1}$ satisfies $\mu_1^n = 1$.
    Setting $Z = Z_0$ and $\lambda = \mu_1$ yields the result.
\end{proof}

\subsection{Non-TI descriptions of TI states}

\begin{corollary}[Non-TI realisation of a TI state]%
    \label{cor:non_ti_to_ti}
    \lean{}
    \notready
    \uses{thm:fundamentalTheorem_injective_chain, def:non_ti_chain}
    Let $\mathbf{A} = (A_0, \dots, A_{n-1})$ be an injective non-TI
    chain of length $n \ge 3$.
    Suppose the state generated by $\mathbf{A}$ is translation-invariant,
    i.e.\ it equals the state generated by the cyclically shifted chain
    $(A_1, \dots, A_{n-1}, A_0)$.
    Then $\mathbf{A}$ admits a translation-invariant description with the
    same bond dimension: there exists an injective tensor $B$ with
    $B^i = Z\, A_0^i\, Z^{-1}$ for some $Z \in \GL_D(\C)$, such that
    the TI chain $(B, \dots, B)$ generates the same state as $\mathbf{A}$.
\end{corollary}

\begin{proof}
    The cyclic shift maps $\mathbf{A}$ to $\sigma(\mathbf{A})$, and both
    generate the same state by hypothesis.
    Apply Theorem~\ref{thm:fundamentalTheorem_injective_chain} to compare
    $\mathbf{A}$ and $\sigma(\mathbf{A})$: there are invertible
    $Z_0, \dots, Z_{n-1}$ with
    $A_{k+1}^i = Z_k A_k^i Z_{k+1}^{-1}$ for all $k$ and $i$.
    Define $B^i = Z_0^{-1} A_0^i Z_1$; then $A_k^i = Z_k^{-1} B^i \cdots$
    unwinds to express each $A_k$ as a gauge conjugate of $A_0$ accumulated
    through the sequence $Z_0, \dots, Z_k$.
    The cyclic compatibility condition closes the gauge loop and ensures
    that the constant chain $(B, \dots, B)$ generates the same state.
\end{proof}

\subsection{Physical symmetries give virtual gauges}

\begin{corollary}[Physical symmetry implies virtual gauge]%
    \label{cor:symmetry_virtual_gauge}
    \lean{}
    \notready
    \uses{cor:ft_ti_chain, def:injective}
    Let $A$ be an injective MPS tensor and let $U \in \GL_d(\C)$ be a
    unitary acting on the physical index.
    If the $n$-site state satisfies
    \[
        U^{\otimes n}\, \ket{\psi(A)} = \ket{\psi(A)}
    \]
    for every $n \ge 3$, then there exist $V \in \GL_D(\C)$ and a phase
    $\varphi \in \C^{\times}$ with $|\varphi| = 1$ such that
    \[
        \sum_{j=0}^{d-1} U_{ij}\, A^j
        = \varphi\, V\, A^i\, V^{-1}
    \]
    for every physical index~$i$.
\end{corollary}

\begin{proof}
    Define $B^i := \sum_j U_{ij} A^j$.
    Since $U$ is invertible, the matrices $\{B^i\}$ span $\MN{D}$, so
    $B$ is injective.
    The symmetry condition $U^{\otimes n} \ket{\psi(A)} = \ket{\psi(A)}$
    means that the TI chain with tensor $B$ generates the same state as
    the TI chain with tensor $A$.
    Apply Corollary~\ref{cor:ft_ti_chain}: there exist $V \in \GL_D(\C)$
    and $\lambda$ with $\lambda^n = 1$ and
    $B^i = \lambda\, V\, A^i\, V^{-1}$ for all~$i$.
    Since this holds for all $n \ge 3$, taking two coprime values of $n$
    forces $\lambda = \varphi$ to be a root of unity that is simultaneously
    an $n$-th root for coprime $n$, hence $|\varphi| = 1$.
\end{proof}

\subsection{Normal MPS (blocked injectivity)}

\begin{corollary}[Normal MPS Fundamental Theorem]%
    \label{cor:normal_mps_ft}
    \lean{}
    \notready
    \uses{thm:fundamentalTheorem_injective_chain, def:injective}
    Let $A$ and $B$ be MPS tensors (not necessarily injective at the
    single-site level) such that the $L$-blocked tensors $A^{[L]}$ and
    $B^{[L]}$ are injective for some blocking length $L \ge 1$.
    If the chains of length $n \ge 3L$ generated by $A$ and $B$
    produce the same state, then $A^{[L]}$ and $B^{[L]}$ are cyclically
    gauge equivalent, i.e.\ there exist invertible matrices
    $Z_0, \dots, Z_{\lfloor n/L \rfloor - 1}$ such that
    \[
        (B^{[L]})^{i_0 \cdots i_{L-1}}_k
        = Z_k\, (A^{[L]})^{i_0 \cdots i_{L-1}}_k\, Z_{k+1}^{-1}
    \]
    for all sites $k$ and physical multi-indices $(i_0, \dots, i_{L-1})$.
\end{corollary}

\begin{proof}
    Blocking every $L$ consecutive sites of the original length-$n$ chain
    produces a new chain of length $m = \lfloor n/L \rfloor \ge 3$
    whose local tensors are $A^{[L]}$ and $B^{[L]}$, which are injective
    by hypothesis.
    The state generated by the blocked chain equals the state generated by
    the original chain (blocking is an isometry on the physical Hilbert
    space).
    Apply Theorem~\ref{thm:fundamentalTheorem_injective_chain} to the
    blocked chains to obtain the cyclic gauge equivalence.
\end{proof}